\chapter{{\fontsize{14pt}{21pt}\selectfont\MakeUppercase{Приложение Д. Листинги ключевого кода}}}
\label{app:listings}

Ниже приведены фрагменты исходного кода программного комплекса (монорепозиторий). Нумерация листингов соответствует приложению~Д.\cite{docker_compose_docs,pydantic_docs,typescript_docs}

\section*{Д.1 Серверная часть}

\noindent\textbf{Листинг Д.1 -- Функция подбора следующего вопроса (\texttt{backend/app/services/adaptive.py})}

\begin{lstlisting}[language=Python]
def select_next_question(db: Session, attempt: Attempt) -> Question | None:
    questions = db.scalars(
        select(Question).where(
            Question.test_id == attempt.test_id,
            Question.tenant_id == attempt.tenant_id,
        )
    ).all()
    remaining = [q for q in questions if q.id not in (attempt.asked_question_ids or [])]
    if not remaining:
        return None

    asked_ids = attempt.asked_question_ids or []
    topic_frequency: Counter[int] = Counter()
    if asked_ids:
        topic_frequency.update(
            db.scalars(
                select(QuestionTopic.topic_id).where(QuestionTopic.question_id.in_(asked_ids))
            )
        )

    question_topic_map: dict[int, list[int]] = defaultdict(list)
    remaining_ids = [question.id for question in remaining]
    if remaining_ids:
        for question_id, topic_id in db.execute(
            select(QuestionTopic.question_id, QuestionTopic.topic_id).where(
                QuestionTopic.question_id.in_(remaining_ids)
            )
        ).all():
            question_topic_map[question_id].append(topic_id)

    def sort_key(question: Question) -> tuple[int, int, int, int]:
        topic_ids = question_topic_map.get(question.id, [])
        if topic_ids:
            min_coverage = min(topic_frequency.get(topic_id, 0) for topic_id in topic_ids)
            total_coverage = sum(topic_frequency.get(topic_id, 0) for topic_id in topic_ids)
        else:
            min_coverage = 0
            total_coverage = 0
        return (
            abs(question.difficulty - attempt.current_difficulty),
            min_coverage,
            total_coverage,
            question.id,
        )

    return sorted(remaining, key=sort_key)[0]
\end{lstlisting}

\noindent\textbf{Листинг Д.2 -- Выдача токенов доступа (\texttt{backend/app/services/auth.py})}

\begin{lstlisting}[language=Python]
def issue_tokens(db: Session, user: User) -> tuple[str, str]:
    access = create_token(str(user.id), "access", settings.access_token_minutes)
    refresh = create_token(str(user.id), "refresh", settings.refresh_token_minutes)
    db.add(
        RefreshToken(
            user_id=user.id,
            token=refresh,
            expires_at=utcnow() + timedelta(minutes=settings.refresh_token_minutes),
            revoked=False,
        )
    )
    return access, refresh
\end{lstlisting}

\noindent\textbf{Листинг Д.3 -- Роли и сущности пользователя (\texttt{backend/app/models/models.py})}

\begin{lstlisting}[language=Python]
class RoleName(StrEnum):
    learner = "learner"
    teacher = "teacher"
    org_admin = "org_admin"
    system_admin = "system_admin"


class Tenant(Base):
    __tablename__ = "tenants"

    id: Mapped[int] = mapped_column(primary_key=True)
    name: Mapped[str] = mapped_column(String(120), unique=True, index=True)
    code: Mapped[str] = mapped_column(String(64), unique=True, index=True)
    locale: Mapped[str] = mapped_column(String(8), default="ru")
    is_active: Mapped[bool] = mapped_column(Boolean, default=True)


class User(Base):
    __tablename__ = "users"

    id: Mapped[int] = mapped_column(primary_key=True)
    email: Mapped[str] = mapped_column(String(255), unique=True, index=True)
    full_name: Mapped[str] = mapped_column(String(255))
    password_hash: Mapped[str] = mapped_column(String(255))
\end{lstlisting}

\noindent\textbf{Листинг Д.4 -- Фрагмент сводной аналитики (\texttt{backend/app/services/analytics.py})}

\begin{lstlisting}[language=Python]
def dashboard_stats(db: Session, tenant_id: int, course_ids: list[int] | None = None) -> dict:
    selected_course_ids = [int(course_id) for course_id in (course_ids or [])]
    has_filter = bool(selected_course_ids)
    enrollment_filters = [Enrollment.tenant_id == tenant_id]
    test_filters = [Test.tenant_id == tenant_id]
    attempt_filters = [Attempt.tenant_id == tenant_id, Attempt.status == "in_progress"]
    course_filters = [Course.tenant_id == tenant_id]
    if has_filter:
        enrollment_filters.append(Enrollment.course_id.in_(selected_course_ids))
        test_filters.append(Test.course_id.in_(selected_course_ids))
        course_filters.append(Course.id.in_(selected_course_ids))
    avg_progress = db.scalar(select(func.avg(Enrollment.progress_percent)).where(*enrollment_filters))
    active_learners = (
        db.scalar(select(func.count(func.distinct(Enrollment.user_id))).where(*enrollment_filters))
        if has_filter
        else db.scalar(select(func.count(Membership.id)).where(Membership.tenant_id == tenant_id))
    )
    return {"users": active_learners or 0, "avg_progress": int(avg_progress or 0)}
\end{lstlisting}

\section*{Д.2 Конфигурация развёртывания}

\noindent\textbf{Листинг Д.5 -- Файл \texttt{docker-compose.yml} (фрагмент сервисов)}

\begin{lstlisting}
services:
  db:
    image: postgres:16-alpine
    ports:
      - "127.0.0.1:5433:5432"
  backend:
    build:
      context: ./backend
    ports:
      - "127.0.0.1:8000:8000"
  web:
    build:
      context: ./web
    ports:
      - "127.0.0.1:8081:80"
\end{lstlisting}

\section*{Д.3 Мобильный клиент}

\noindent\textbf{Листинг Д.6 -- Экран прохождения теста (\texttt{mobile/lib/main.dart}, фрагмент)}

\begin{lstlisting}[language=Java]
class TestFlowScreen extends StatefulWidget {
  const TestFlowScreen({
    super.key,
    required this.api,
    required this.session,
    required this.testId,
  });

  final ApiClient api;
  final SessionState session;
  final int testId;

  @override
  State<TestFlowScreen> createState() => _TestFlowScreenState();
}

class _TestFlowScreenState extends State<TestFlowScreen> {
  Future<void> start() async {
    final started =
        await widget.api.startAttempt(widget.testId, widget.session);
    attemptId = started['attempt_id'] as int;
    await loadNextQuestion();
  }

  Future<void> loadNextQuestion() async {
    final nextQuestion =
        await widget.api.getAttemptQuestion(attemptId!, widget.session);
    if (nextQuestion == null) {
      await finishAttempt();
      return;
    }
    setState(() {
      question = nextQuestion;
      questionShownAt = DateTime.now();
    });
  }
}
\end{lstlisting}

\textit{Примечание: пароли, секретные ключи и продакшен-URL в листингах не приводятся; при локальной отладке используются значения из \texttt{.env.example}.}
